\section{Architektur-Patterns}

\subsection{Load Balancer}
Die Hauptaufgabe eines \textbf{Load Balancers} ist der \emph{Lastenausgleich}. Wenn viele Nutzer Anfragen an eine Anwendung stellen, werden diese auf mehrere Server verteilt. Dadurch werden Über- und Unterlastung von Ressourcen vermieden. Die meisten PaaS-Angebote in der Cloud kommen mit einem vorkonfigurierten Load Balancer. Bei einem Load Balancer skaliert man \emph{horizontal} (mehr Instanzen, nicht stärkere Hardware).

Weitere Funktionen eines Load Balancers:
\begin{itemize}
    \item \textbf{Health Checks:} Erkennt ausgefallene Instanzen und leitet Traffic automatisch um
    \item \textbf{SSL-Terminierung:} Verschlüsselung wird zentral am Load Balancer beendet
    \item \textbf{Session Persistence (Sticky Sessions):} Nutzer werden immer zum selben Server geleitet
    \item \textbf{Ressource-Cluster:} Zusammenfassung mehrerer Instanzen zu einem logischen Pool
    \item \textbf{Netzwerkisolierung:} Backend-Server sind nicht direkt aus dem Internet erreichbar
\end{itemize}

\subsection{CDN - Content Delivery Network}
Ein Content Delivery Network ist ein weltweit verteiltes Netz aus Cache-Servern,
das Inhalte möglichst nah am Endnutzer ausliefert. Während ein Load Balancer den
Traffic \emph{innerhalb} einer Region verteilt, verkürzt ein CDN die \textbf{Latenz}
geografisch, indem es Inhalte an vielen Standorten zwischenspeichert. Dadurch wird
gleichzeitig der \textbf{Origin-Server} (Ursprungsserver) entlastet, da wiederholte
Anfragen direkt aus dem Cache beantwortet werden.

Zentrale Begriffe:
\begin{itemize}
    \item \textbf{Edge Location / PoP (Point of Presence):} Standort eines
          CDN-Cache-Servers nah am Nutzer. Liegt noch \emph{unterhalb} der Region
          (vgl. Kapitel \enquote{Regionen und Verfügbarkeitszonen}) und damit
          deutlich näher am Endgerät.
    \item \textbf{Origin:} Der Ursprungsserver, auf dem die Inhalte tatsächlich
          liegen (z.\,B. ein Object Storage oder eine Web-App).
    \item \textbf{Cache-Hit / Cache-Miss:} Bei einem Hit liefert die Edge Location
          den Inhalt direkt aus; bei einem Miss holt sie ihn zuerst vom Origin und
          legt ihn dann lokal ab.
    \item \textbf{TTL (Time to Live):} Legt fest, wie lange ein Inhalt im Cache
          gültig bleibt, bevor er erneut vom Origin geladen wird.
\end{itemize}

CDNs eignen sich vor allem für \textbf{statische Inhalte} (Bilder, Videos,
CSS/JS, Downloads). Für \textbf{dynamische Inhalte} ist der Nutzen geringer, da
diese pro Anfrage individuell erzeugt werden. Beispiele für CDN-Dienste sind
\textit{Azure Front Door / Azure CDN}, \textit{Amazon CloudFront} und
\textit{Cloudflare}.

\subsection{Cloud Bursting}
Wenn die \textit{On-Premises}-Ressourcen ausgelastet sind, werden zusätzliche Workloads \emph{dynamisch} in die Cloud ausgelagert. Sobald die Last sinkt, kehren die Workloads wieder auf die eigene Infrastruktur zurück. So zahlt man Cloud-Ressourcen \emph{nur bei tatsächlichem Bedarf}.

\subsection{Disaster Recovery}
\textbf{Disaster Recovery (DR)} bezeichnet Strategien und Maßnahmen, mit denen eine Anwendung oder ein System nach einem \emph{schwerwiegenden Ausfall} (z.\,B. Rechenzentrumsausfall, Datenverlust, Cyberangriff) wiederhergestellt werden kann.

\subsubsection{Schlüsselkennzahlen}
\begin{itemize}
    \item \textbf{RTO (Recovery Time Objective):} Maximale tolerierbare Ausfallzeit -- wie lange darf die Wiederherstellung dauern?
    \item \textbf{RPO (Recovery Point Objective):} Maximaler tolerierbarer Datenverlust -- wie alt darf das letzte Backup sein?
\end{itemize}
Ein niedriges RTO/RPO bedeutet höhere Kosten, da mehr Redundanz und häufigere Replikation nötig sind.

\subsubsection{DR-Strategien (nach Kosten/Geschwindigkeit)}
\begin{itemize}
    \item \textbf{Backup \& Restore:} Günstigste Variante; Daten werden regelmäßig gesichert und im Ernstfall wiederhergestellt. Hohes RTO/RPO.
    \item \textbf{Pilot Light:} Minimale Kerninfrastruktur läuft dauerhaft in der DR-Region; bei Bedarf wird sie hochskaliert.
    \item \textbf{Warm Standby:} Vollständige, aber kleiner dimensionierte Kopie des Systems läuft parallel; wird im Ernstfall hochskaliert.
    \item \textbf{Active/Active (Multi-Site):} Zwei oder mehr vollständige Systeme laufen gleichzeitig aktiv; kein Failover nötig, maximale Verfügbarkeit. Teuerste Variante.
\end{itemize}

\subsection*{Azure Storage-Redundanzoptionen}
Azure bietet verschiedene Redundanzstufen für Speicher, die sich nach geografischer Verteilung und Replikationsart unterscheiden:

\begin{itemize}
    \item \textbf{LRS} (Locally Redundant Storage): 3 synchrone Kopien innerhalb \emph{eines} Rechenzentrums in einer Region. Günstigste Option, schützt nur vor Hardware-Ausfall.
    \item \textbf{ZRS} (Zone-Redundant Storage): 3 synchrone Kopien in 3 verschiedenen \emph{Verfügbarkeitszonen} innerhalb einer Region. Schützt vor dem Ausfall einer Zone.
    \item \textbf{GRS} (Geo-Redundant Storage): LRS in der primären Region \emph{plus} asynchrone Replikation in eine sekundäre, geografisch getrennte Partnerregion (dort ebenfalls als LRS gespeichert). Insgesamt 6 Kopien. Schützt vor einem regionalen Ausfall; Lesezugriff auf die Sekundärregion nur nach einem Failover.
    \item \textbf{RA-GRS} (Read-Access GRS): Wie GRS, aber mit dauerhaftem Lesezugriff auf die Sekundärregion ohne Failover.
    \item \textbf{GZRS} (Geo-Zone-Redundant Storage): ZRS in der primären Region \emph{plus} asynchrone Replikation in eine sekundäre Region (dort als LRS). Höchste Redundanz -- schützt vor Zonen- \emph{und} Regionsausfall.
    \item \textbf{RA-GZRS} (Read-Access GZRS): Wie GZRS, aber mit dauerhaftem Lesezugriff auf die Sekundärregion.
\end{itemize}

\subsection{Service Discovery}
\textbf{Service Discovery} bedeutet die \emph{automatische Erkennung von Diensten} in einem Netzwerk. Services \emph{registrieren sich selbst}, sodass andere Services sie automatisch finden, ohne dass IP-Adressen oder Ports hart kodiert werden müssen. Durch Service Discovery kann man leicht Services anbinden. Gerade in \textit{Multi-Cloud-Umgebungen}, bei denen ein virtuelles Netzwerk zwischen privater und öffentlicher Cloud hergestellt wird, ist die automatische Erkennung von Services besonders wichtig.

\subsection{VPC – Virtual Private Cloud}
Eine \textbf{Virtual Private Cloud (VPC)} ist ein \emph{logisch isoliertes} Netzwerksegment
innerhalb einer öffentlichen Cloud-Infrastruktur. Obwohl die physische Hardware
mit anderen Kunden geteilt wird, verhält sich eine VPC wie ein eigenes privates
Netzwerk: mit eigenen IP-Bereichen, Subnetzen, Routing-Regeln und Firewalls.
Der Zugriff von außen wird über klar definierte Gateways und \textit{Security Groups}
kontrolliert.

\subsection{Allocating -- Ressourcenzuweisung und Kostenzuordnung}
\textbf{Allocating} bezeichnet die gezielte Zuweisung von Cloud-Ressourcen zu bestimmten
Projekten, Abteilungen oder Workloads. Ziel ist es, den Ressourcenverbrauch
nachvollziehbar zu machen und Kosten \emph{verursachergerecht} zuzuordnen.
In der Praxis wird dies über sogenannte \textit{Tags} (Key-Value-Pair)
realisiert, die jeder Ressource beim Erstellen mitgegeben werden
Cloud-Portale wie das \textit{Azure Cost Management} nutzen diese Tags, um automatisierte
Kostenberichte und Budgetalarmierungen zu erstellen.
\subsection{Provisioning -- Dynamische Bereitstellung}
\textbf{Provisioning} beschreibt den Prozess, Cloud-Ressourcen bereitzustellen und in
einen nutzbaren Zustand zu versetzen. Es gibt dabei \emph{drei grundlegende Ansätze}:

\begin{itemize}
    % Manuelles Provisioning
    \item \textbf{Manuelles Provisioning:} Ressourcen werden direkt über
    Cloud-Portale wie \texttt{portal.azure.com} oder die AWS-Management-Console
    angelegt. Dieser Ansatz ist einfach, aber fehleranfällig und schlecht
    reproduzierbar.

    % IaC
    \item \textbf{IaC (Infrastructure as Code):} Ressourcen werden
    deklarativ als Code beschrieben und versioniert. Beispiele sind
    \textit{Bicep} und \textit{ARM-Templates} für Azure sowie
    \textit{Terraform} als Cloud-agnostische Alternative. Der Vorteil liegt in
    der Reproduzierbarkeit und der Integration in CICD-Pipelines.

    % Self-Service
    \item \textbf{Self-Service-Provisioning:} Entwickler können aus einem
    vordefinierten Katalog genehmigter Ressourcen selbst auswählen und diese
    eigenständig bereitstellen. Damit werden Wartezeiten auf das Ops-Team
    reduziert, während Governance-Regeln trotzdem eingehalten werden.
\end{itemize}

\subsection{Virtual Server Clustering}
Beim Virtual Server Clustering werden mehrere virtuelle Maschinen zu einer
logischen Einheit zusammengefasst, die gemeinsam Anfragen verarbeiten.
Ziel ist die Steigerung von Verfügbarkeit und Ausfallsicherheit.

Man unterscheidet dabei zwei grundlegende Betriebsmodi:

\begin{itemize}
    % Active-Passive
    \item \textbf{Active-Passive-Cluster:} Ein Knoten ist aktiv und verarbeitet
    Anfragen, ein zweiter Knoten steht im Standby. Bei einem Ausfall übernimmt
    der Standby-Knoten automatisch (\textit{Failover}).

    % Active-Active
    \item \textbf{Active-Active-Cluster:} Alle Knoten verarbeiten gleichzeitig
    Anfragen. Ein vorgelagerter \textit{Load Balancer} verteilt den eingehenden
    Traffic gleichmäßig. Fällt ein Knoten aus, übernehmen die restlichen
    Knoten dessen Last.
\end{itemize}

In modernen Cloud-Umgebungen werden Cluster häufig über
\textit{Virtual Machine Scale Sets} (Azure) oder
\textit{Auto Scaling Groups} (AWS) automatisch verwaltet.

\subsection{Zero Downtime}
Das Ziel von Zero-Downtime-Deployments ist es, neue Softwareversionen
auszuliefern, ohne dass Endnutzer eine Unterbrechung wahrnehmen. Typische
Strategien sind:

\begin{itemize}
    % Blue-Green
    \item \textbf{Blue-Green-Deployment:} Es existieren zwei identische
    Produktionsumgebungen (Blau und Grün). Während Blau produktiv läuft, wird
    die neue Version in Grün deployt und getestet. Anschließend wird der
    Traffic per Load Balancer auf Grün umgeschaltet. Bei Problemen ist ein
    sofortiges Zurückschalten möglich (\textit{Rollback}).

    % Rolling Update
    \item \textbf{Rolling Update:} Instanzen einer Anwendung werden
    schrittweise ausgetauscht. Zu jedem Zeitpunkt laufen alte und neue
    Versionen parallel, bis alle Instanzen aktualisiert sind. Kubernetes nutzt
    dieses Verfahren standardmäßig.

    % Canary Release
    \item \textbf{Canary Release:} Die neue Version wird zunächst nur an einen
    kleinen Prozentsatz der Nutzer ausgeliefert. Verhält sie sich stabil, wird
    der Rollout schrittweise auf alle Nutzer ausgeweitet.
\end{itemize}


\subsection{VM Migration}
VM Migration bezeichnet das Verschieben einer virtuellen Maschine von einem
Host oder Standort zu einem anderen. Man unterscheidet zwei Hauptvarianten:

\begin{itemize}
    % Live Migration
    \item \textbf{Live Migration (Hot Migration):} Die VM wird im
    laufenden Betrieb auf einen anderen physischen Host verschoben. Der
    Arbeitsspeicher wird dabei inkrementell übertragen, sodass die Ausfallzeit
    gegen null geht. VMware bezeichnet dieses Verfahren als \textit{vMotion},
    Hyper-V als \textit{Live Migration}.

    % Cold Migration
    \item \textbf{Cold Migration:} Die VM wird heruntergefahren,
    verschoben und anschließend neu gestartet. Sie ist einfacher umzusetzen,
    verursacht aber eine messbare Downtime.

    \item \textbf{Warm Migration:} Eine Mischform aus beiden Verfahren.
    Im \textit{Precopy-Modus} wird der Großteil der Daten bereits bei
    laufender VM kopiert. Erst die verbleibenden Restdaten werden in einer
    kurzen \textit{Cutover-Phase} nach dem Stoppen der VM übertragen. Dadurch
    ergibt sich eine deutlich reduzierte Downtime gegenüber der Cold Migration.
\end{itemize}

Im Cloud-Kontext wird VM Migration auch für die Verlagerung von
On-Premises-Workloads in die Cloud genutzt (\textit{Lift and Shift}).
Azure stellt dafür den Dienst \textit{Azure Migrate} bereit, der Abhängigkeiten
analysiert und den Migrationsprozess begleitet.
